\documentclass[a4paper,10.5pt]{article}
\usepackage[utf8]{inputenc}
\usepackage[T1]{fontenc}
\usepackage[french]{babel}
\usepackage[left=1cm,right=1cm,top=.5cm,bottom=0cm]{geometry}
\usepackage{enumitem}
\usepackage{hyperref}
\usepackage{xcolor}
\usepackage{titlesec}
\usepackage{fontawesome5}
\usepackage{lmodern}
\usepackage[normalem]{ulem}

% --- Couleurs ---
\definecolor{primary}{RGB}{35, 55, 123}
\definecolor{secondary}{RGB}{80, 80, 80}

% --- Configuration des liens ---
\hypersetup{
    colorlinks=true,
    linkcolor=primary,
    filecolor=primary,
    urlcolor=primary,
}

% --- Formatage des titres ---
\titleformat{\section}{\large\bfseries\color{primary}\uppercase}{}{0em}{}[\titlerule]
\titlespacing{\section}{0pt}{8pt}{5pt}

% --- Commandes personnalisées ---
\newcommand{\resumeItem}[1]{\item\small{#1}}
\newcommand{\resumeSubheading}[4]{
  \vspace{-1pt}\item
    \begin{tabular*}{0.97\textwidth}[t]{l@{\extracolsep{\fill}}r}
      \textbf{#1} & \textit{\small #2} \\
      \textit{\small #3} & \textit{\small #4} \\
    \end{tabular*}\vspace{-5pt}
}

\begin{document}

% --- EN-TÊTE ---
\begin{center}
    {\Large \textbf{Loïc Aron MBASSI EWOLO}} \\ \vspace{4pt}
    \textbf{\large Stage Ingénieur Structuration de Modèles \& Simulation} \\
    \textit{\small Disponible dès \textbf{Avril 2026} pour 4 à 6 mois} \\ \vspace{4pt}
    \small
    \faMapMarker\ Cergy, France (Mobile) \quad
    \faPhone\ (+33) 7 59 52 33 98 \quad
    \faEnvelope\ \href{mailto:loicaron.mbassiewolo@ensea.fr}{loicaron.mbassiewolo@ensea.fr} \\ \vspace{2pt}
    \faLinkedin\ \href{https://www.linkedin.com/in/loic-aron-mbassi-ewolo-3942a9246/}{linkedin} \quad
    \faGithub\ \href{https://github.com/Nameless0l}{github} \quad
    \faGlobe\ \href{https://mbassiloic.tech}{mbassiloic.tech}
\end{center}

\vspace{-6pt}

% --- PROFIL ---
\noindent\small{Élève-ingénieur en double diplôme \textbf{ENSEA/Polytechnique Yaoundé}, spécialisé en \textbf{modélisation UML}, développement \textbf{Python/C++} et architectures de systèmes complexes. Créateur d'un outil open source de génération de code à partir de diagrammes UML. Je recherche un stage chez \textbf{Framatome} pour appliquer mes compétences en ingénierie basée sur les modèles (MBSE) et simulation de systèmes.}

% --- FORMATION ---
\section{Formation}
\begin{itemize}[leftmargin=*, label=\tiny$\bullet$, nosep]
    \item \textbf{ENSEA (École Nationale Supérieure de l'Électronique et de ses Applications)} \hfill \textit{Cergy, France | 2025 -- Présent} \\
    \small{Ingénieur 2A, Double Diplôme - Systèmes Embarqués, Traitement du Signal et IA.} \\
    \textit{\small{Cours : Conception numérique, Modélisation systèmes, Programmation C/C++.}}
    \vspace{2pt}
    \item \textbf{École Polytechnique de Yaoundé (1\textsuperscript{ère} école d'ingénieur CEMAC)} \hfill \textit{Cameroun | 2021 -- 2025} \\
    \small{Diplôme d'Ingénieur de Conception (Master 2) : \textbf{Modélisation UML/AUML}, Génie Logiciel, Algorithmique.}
    \item \textbf{Université de Yaoundé I} \hfill \textit{Cameroun | 2020 -- 2022} \\
    \small{Licence Informatique \& Mathématiques : Algorithmique C, Analyse, Algèbre Linéaire.}
\end{itemize}

% --- COMPÉTENCES TECHNIQUES ---
\section{Compétences Techniques}
\begin{itemize}[leftmargin=*, nosep]
    \item \small{\textbf{Langages :} \textbf{Python} (Avancé), \textbf{C/C++}, Java, PHP, JavaScript/TypeScript.}
    \item \small{\textbf{Modélisation \& Architecture :} \textbf{UML/AUML} (Diagrammes de classes, séquence, activité), MBSE, Design Patterns.}
    \item \small{\textbf{Systèmes \& Simulation :} Linux, Programmation concurrente (IPC, Signaux), Docker, Algorithmes d'optimisation.}
    \item \small{\textbf{Outils :} Git, Draw.io, XML Parsing, CI/CD, Architecture Microservices.}
\end{itemize}

% --- PROJETS CLÉS ---
\section{Projets Académiques \& Personnels}
\begin{itemize}[leftmargin=*, label={-}]

\item \textbf{Laravel API Generator - Génération de Code depuis UML/XML} \hfill \textit{Projet Open Source | 2024}
\vspace{-2pt}
\begin{itemize}[leftmargin=0.4cm, nosep, label=\tiny$\bullet$]
    \item \small{Développement d'un outil \textbf{Python} générant automatiquement une architecture backend complète depuis des \textbf{diagrammes UML}.}
    \item \small{\textbf{Parsing XML} avancé : analyse géométrique des liens entre entités exportés par Draw.io.}
    \item \small{Génération automatique de modèles, contrôleurs, migrations et routes à partir de la structure UML.}
    \item \small{\textbf{+30 étoiles GitHub}, déployé sur Packagist avec +20 téléchargements.}
\end{itemize}
\vspace{3pt}

\item \textbf{Projet TAXI - Simulation Système Temps Réel Multi-Processus} \hfill \textit{Projet Académique | 2024}
\vspace{-2pt}
\begin{itemize}[leftmargin=0.4cm, nosep, label=\tiny$\bullet$]
    \item \small{Développement en \textbf{C} d'un noyau de simulation pour une flotte de taxis avec architecture modulaire.}
    \item \small{Gestion de la \textbf{mémoire partagée} (SHM), communication inter-processus (IPC) et synchronisation par signaux UNIX.}
    \item \small{Implémentation d'un ordonnanceur FIFO avec prévention des interblocages.}
    \item \small{Visualisation temps réel avec \textbf{OpenGL}.}
\end{itemize}
\vspace{3pt}

\item \textbf{Système Multi-Agents Agricole - Orchestration \& Simulation} \hfill \textit{Projet Personnel | 2024-2025}
\vspace{-2pt}
\begin{itemize}[leftmargin=0.4cm, nosep, label=\tiny$\bullet$]
    \item \small{Conception d'un système orchestrant \textbf{5 agents IA} spécialisés avec résolution de conflits.}
    \item \small{Architecture modulaire avec coordination automatique inter-agents et adaptation contextuelle (RAG).}
    \item \small{Technologies : Google ADK, \textbf{Python} (Poetry, FastAPI), Docker.}
\end{itemize}
\vspace{3pt}

\item \textbf{GoFind - Architecture Microservices} \hfill \textit{Projet Académique | 2024}
\vspace{-2pt}
\begin{itemize}[leftmargin=0.4cm, nosep, label=\tiny$\bullet$]
    \item \small{Conception d'une architecture \textbf{microservices} complète avec communication asynchrone (RabbitMQ).}
    \item \small{Modélisation des flux de données et \textbf{diagrammes d'architecture} entre services.}
    \item \small{Technologies : Laravel, NestJS, MongoDB, Docker.}
\end{itemize}
\vspace{3pt}

\item \textbf{Upgrade Jetson - Robotique Autonome (Défense)} \hfill \textit{Challenge CoHoMa | En cours}
\vspace{-2pt}
\begin{itemize}[leftmargin=0.4cm, nosep, label=\tiny$\bullet$]
    \item \small{Optimisation d'algorithmes de cartographie (SLAM) et pilotage pour robots militaires autonomes.}
    \item \small{Intégration et \textbf{simulation} de modèles IA embarqués sur Nvidia Jetson.}
\end{itemize}

\end{itemize}

% --- EXPÉRIENCE PROFESSIONNELLE ---
\section{Expériences Professionnelles}
\begin{itemize}[leftmargin=*, label={-}]

    \resumeSubheading
      {Assistant de Recherche - Généralisation IA (Grokking)}{Juin 2025 -- Sept. 2025}
      {MILA (Institut Québécois d'IA) - Remote}{Québec, Canada}
      \begin{itemize}[leftmargin=0.4cm, nosep, label=\tiny$\bullet$]
        \item \small{Implémentation de pipelines de tests statistiques sous \textbf{PyTorch/Python}.}
        \item \small{Reproduction d'expériences SOTA et analyse mathématique de la convergence.}
      \end{itemize}
    \vspace{2pt}

    \resumeSubheading
      {Développeur Backend - API Financière}{Oct. 2024 -- Jan. 2025}
      {CSC (Cameroon Software Company) - Fintech}{Yaoundé, Cameroun}
      \begin{itemize}[leftmargin=0.4cm, nosep, label=\tiny$\bullet$]
        \item \small{Conception d'une architecture backend sécurisée avec gestion des accès concurrents.}
        \item \small{Mise en place de tests unitaires et revues de code pour garantir la fiabilité.}
      \end{itemize}

\end{itemize}

% --- CERTIFICATIONS & SOFT SKILLS ---
\section{Certifications, Leadership \& Langues}
\begin{itemize}[leftmargin=*, nosep]
    \item \textbf{Leadership :} Chef de la Cellule Projets \& IA (Polytechnique Yaoundé), animation de workshops techniques.
    \item \textbf{Hackathons :} \textbf{1\textsuperscript{ère} Place} IndabaX Africa (IA Santé), \textbf{1\textsuperscript{ère} Place} CITS National.
    \item \textbf{Certifications :} Supervised ML (DeepLearning.AI), Python for Data Science (IBM).
    \item \textbf{Langues :} Français (Natif), \textbf{Anglais} (Professionnel/Technique - B2).
    \item \textbf{Intérêts :} Piano, Rédaction d'articles techniques, Veille technologique.
\end{itemize}

\end{document}
